\documentclass{beamer}
\usefonttheme{professionalfonts}% Now beamer didn't modify the math fonts
\usepackage{geometry}
\input{macros.tex}
% \usepackage{tikz-cd}
\beamertemplatenavigationsymbolsempty
\setbeamertemplate{footline}[page number]{}
\setbeamerfont{footline}{size=\footnotesize}
\usepackage{tcolorbox}
\usepackage{biblatex}
% \usepackage{tikzcd}
\usetheme{}
\usepackage{tikz}

\definecolor{asparagus}{rgb}{0.53, 0.66, 0.42}


\title[]{Clase tutorial 11: Simulación de Monte Carlo.}
% \subtitle{Workshop de LOREL 2024.}
\date{30 de Octubre de 2024.}
\author[]{ Investigación operativa.}
\institute{Universidad de San Andrés}
% \titlegraphic{\hfill\includegraphics[height=1.5cm]{logo.pdf}}

\begin{document}
\maketitle

\begin{frame}{Objetivos de la clase de hoy.}
  La idea de la clase de hoy es ver los siguientes temas:
  \begin{itemize}
    \item Simulación de Monte Carlo.
  \end{itemize}
\end{frame}

\begin{frame}[fragile]{Problema 1.}
En un semáforo se acumulan autos cada vez que su luz se pone roja.
Supongamos que cada vez que el semáforo se pone en verde todos los autos
acumulados avanzan. 
Supongamos también que la cantidad de autos a acumularse en cada iteración es
independiente de la historia y de los estados anteriores.
Si la cantidad de autos a acumularse en cada luz roja sigue la siguiente distribución de probabilidad.
\begin{center}
  \begin{adjustbox}{width = 0.85 \textwidth}
      \begin{tabular}{|c|c|c|c|c|c|}
        \hline
        Cantidad de Autos & 0    & 1   & 2    & 3    & 4 \\
        \hline 
        $\mathbb{P}$      & 0.15 & 0.2 & 0.35 & 0.23 & 0.07 \\
        \hline 
      \end{tabular}
  \end{adjustbox}
\end{center}

Simular 10, 100 y 1000 veces el semáforo en rojo.

\end{frame}


\begin{frame}[fragile]{Problema 2.}
 Utilizar el método de simulación de Monte Carlo para aproximar la constante $\pi$.
\end{frame}

\begin{frame}[fragile]{Problema 3.}
  El \textbf{craps} es un juego de dados sencillo dado por las siguientes reglas.
  Tiramos dos dados.
  Si suman 7 o un 11 entonces ganamos de inmediato.
  Si suman 2,3 o 12 entonces perdemos de inmeadiato.
  En cualquier otro caso volvemos a tirar.

  \begin{enumerate}[a)]
    \item Pensar analíticamente qué porcentaje de veces deberíamos ganar.

    \item Resolverlo empíricamente por medio de simulaciones.
  \end{enumerate}
\end{frame}


\begin{frame}[fragile]{Tarea.}
  
La \textbf{guerra} es un juego sencillo de cartas que se juega de a dos personas.
A cada persona se le reparte la mitad de un mazo de 52 cartas francesas. 
En cada turno ambos jugadores tiran la primera carta de su mazo y la comparan. 
Quien tenga la carta más alta gana un punto y en caso que ambas cartas sean iguales entonces la próxima carta decide quién se queda con dos puntos y así continua el juego hasta que se agota el mazo. Eventualmente el ganador es quien consiga más puntos.





\begin{enumerate}[a)]
  \item Simular una partida entre dos manos arbitrarias y anotarlo en una tabla (no hace falta hacer un programa que nos de la tabla.)
\end{enumerate}
\end{frame}

\end{document}
